% ===========================================================================
%  Chapitre 6 — Étude complète d'une fonction
%  PE 2026, composante ANALYSE — RAPE, deuxième période
% ===========================================================================
\bmtchapitre{Étude complète d'une fonction}
  {Mener l'étude complète d'une fonction rationnelle : ensemble de
   définition, limites, dérivée, tableau de variation, points particuliers,
   position de la courbe par rapport à une droite, et tracé dans un repère.}
  {Analyse \textbullet\ environ 5 heures}

Ce chapitre n'apporte presque aucune notion nouvelle : il assemble. Les
limites viennent du chapitre 4, la dérivée du chapitre 5, le tableau de
signes du chapitre 1. Ce qui s'apprend ici, c'est \emph{l'ordre} dans lequel
on les met — et cet ordre est toujours le même, quelle que soit la fonction.

C'est aussi le chapitre qui donne sa forme au problème du baccalauréat. Que
la fonction contienne un logarithme, une exponentielle ou seulement un
quotient de polynômes, les questions se suivent dans cet ordre-là, session
après session. Apprenez le plan une fois, il servira toute l'année.

% ---------------------------------------------------------------------------
\begin{revision}
Tout vient des chapitres 1, 4 et 5.

\begin{enumerate}[label=\textbf{\arabic*.}, leftmargin=*, itemsep=0.35em]
  \item Donner l'ensemble de définition de $x \mapsto \dfrac{x+1}{x^{2}-9}$.
  \item Calculer $\limpinf \dfrac{2x^{2}+1}{x^{2}-x}$.
  \item Dériver $x \mapsto \dfrac{x^{2}}{x+1}$.
  \item Étudier le signe de $x^{2}-2x-3$.
  \item Écrire l'équation de la tangente en $0$ à la courbe de
        $x \mapsto x^{2}+x+1$.
  \item Que signifie « la courbe est au-dessus de la droite $(D)$ » en
        termes d'inégalité ?
\end{enumerate}
\end{revision}

\begin{automatismes}
Sans calculatrice.

\begin{enumerate}[label=\textbf{\alph*)}, leftmargin=*, itemsep=0.25em]
  \item $\dfrac{4}{0^{+}}$
  \item $\dfrac{-3}{0^{-}}$
  \item $\limpinf \dfrac{x^{2}}{x}$
  \item Signe de $(x-1)^{2}$
  \item Valeur interdite de $\dfrac{5}{2x-6}$
  \item $f(x) = \dfrac{x^{2}}{x-1}$ : calculer $f(2)$
  \item $f(x) = \dfrac{x^{2}}{x-1}$ : calculer $f(0)$
  \item Une courbe coupe $(Oy)$ en un point d'abscisse
\end{enumerate}
\end{automatismes}

\begin{corrige}[automatismes]
a) $+\infty$ \quad b) $+\infty$ \quad c) $+\infty$ \quad d) positif ou nul
\quad e) $x = 3$ \quad f) $4$ \quad g) $0$ \quad h) $0$.
\end{corrige}

% ===========================================================================
\section{Le plan d'étude, une fois pour toutes}

\begin{methode}{étudier une fonction, en sept étapes}
\begin{enumerate}[label=\textbf{\arabic*.}, leftmargin=*, itemsep=0.18em]
  \item \emph{Ensemble de définition.} Écrire $\Df$ et le justifier — un
        dénominateur non nul, une racine carrée positive.
  \item \emph{Parité éventuelle.} Si $\Df$ est symétrique, tester $f(-x)$ ;
        sinon, passer.
  \item \emph{Limites aux bornes de $\Df$.} Une par borne, sans en oublier
        aucune, et chacune suivie de son interprétation graphique.
  \item \emph{Dérivée.} Calculer $f'(x)$, la réduire, factoriser le
        numérateur.
  \item \emph{Signe de $f'$ et tableau de variation.} Ligne de signes,
        flèches, valeurs des extremums, limites aux extrémités.
  \item \emph{Points particuliers.} Intersections avec les axes, tangente en
        un point donné, position par rapport à une droite ou à une
        asymptote.
  \item \emph{Tracé.} Repère, unité, asymptotes en pointillé d'abord, puis
        la courbe.
\end{enumerate}
\end{methode}

\begin{remarque}
Ce plan est celui du sujet d'examen, dans cet ordre. Un candidat qui
connaît les sept étapes sait, en lisant l'énoncé, où il en est et combien de
points restent à prendre.
\end{remarque}

% ===========================================================================
\section{Premier type : \texorpdfstring{$f(x) = \dfrac{ax^{2}+bx+c}{dx+e}$}{quotient de degre 2 sur degre 1}}

Ces fonctions ont une asymptote verticale et une asymptote oblique. La forme
$f(x) = \alpha x+\beta+\dfrac{k}{dx+e}$ donne tout d'un coup.

\begin{exemple}[une étude complète, du début au tracé]
Soit $f(x) = \dfrac{x^{2}-3x+4}{x-2}$.

\emph{1. Ensemble de définition.} $x-2 \neq 0$, donc
$\Df = \R\setminus\{2\}$.

\emph{2. Parité.} $\Df$ n'est pas symétrique par rapport à $0$ : la question
ne se pose pas.

\emph{3. Limites.} On vérifie que $f(x) = x-1+\dfrac{2}{x-2}$ :
\[
  (x-1)(x-2)+2 = x^{2}-3x+2+2 = x^{2}-3x+4 . \checkmark
\]
Donc $f(x) \to \pm\infty$ en $\pm\infty$, et la droite $(D) : y = x-1$ est
asymptote oblique. En $2$, le terme $\frac{2}{x-2}$ tend vers $-\infty$ à
gauche et $+\infty$ à droite : la droite $x = 2$ est asymptote verticale.

\emph{4. Dérivée.}
\[
  f'(x) = \frac{(2x-3)(x-2)-\left(x^{2}-3x+4\right)}{(x-2)^{2}}
  = \frac{x^{2}-4x+2}{(x-2)^{2}} .
\]

\emph{5. Signe et variations.} Le numérateur a pour discriminant
$16-8 = 8$ et pour racines $2-\sqrt2 \approx 0{,}59$ et
$2+\sqrt2 \approx 3{,}41$. Le dénominateur est positif :

\begin{center}\footnotesize
\renewcommand{\arraystretch}{1.3}
\begin{tabular}{|c|ccccccccc|}
\hline
$x$ & $-\infty$ & & $2-\sqrt2$ & & $2$ & & $2+\sqrt2$ & & $+\infty$\\
\hline
$f'(x)$ & & $+$ & $0$ & $-$ & \textbardbl & $-$ & $0$ & $+$ & \\
\hline
$f$ & $-\infty$ & $\nearrow$ & $-1{,}83$ & $\searrow$ & \textbardbl &
$\searrow$ & $4{,}83$ & $\nearrow$ & $+\infty$\\
\hline
\end{tabular}
\end{center}

\emph{6. Points particuliers.} La courbe coupe $(Oy)$ en
$\left(0\,;\,-2\right)$. Elle ne coupe pas $(Ox)$ : $x^{2}-3x+4$ a un
discriminant négatif. Position par rapport à $(D)$ : le signe de
$f(x)-(x-1) = \frac{2}{x-2}$ est celui de $x-2$, donc la courbe est
au-dessus de $(D)$ pour $x>2$ et en dessous pour $x<2$.

\emph{7. Tracé.}

\begin{center}
\begin{tikzpicture}
\begin{axis}[
  width = 11.2cm, height = 7cm,
  axis lines = middle, xlabel = {$x$}, ylabel = {$y$},
  xmin = -3.4, xmax = 7.4, ymin = -6.4, ymax = 9.4,
  xtick = {-3,-2,-1,1,2,3,4,5,6,7}, ytick = {-6,-4,-2,2,4,6,8},
  axis line style = {bmtGris},
  tick label style = {font=\scriptsize, color=bmtGris},
  grid = both, grid style = {bmtGrisClair, line width=0.3pt},
  clip = true, restrict y to domain = -14:16,
]
  \addplot[bmtIndigo, line width=1.4pt, domain=-3.3:1.86, samples=220]
    {(x^2-3*x+4)/(x-2)};
  \addplot[bmtIndigo, line width=1.4pt, domain=2.14:7.3, samples=220]
    {(x^2-3*x+4)/(x-2)};
  \addplot[bmtLaterite, line width=1pt, dash pattern=on 3pt off 2pt,
           domain=-3.3:7.3] {x-1};
  \draw[bmtLaterite, line width=1pt, dash pattern=on 3pt off 2pt]
    (axis cs:2,-6.3) -- (axis cs:2,9.3);
  \filldraw[bmtVetiver] (axis cs:0.586,-1.828) circle (2pt);
  \filldraw[bmtVetiver] (axis cs:3.414,4.828) circle (2pt);
  \node[bmtLaterite, font=\scriptsize, anchor=west]
    at (axis cs:5.6,4.2) {$(D) : y = x-1$};
  \node[bmtLaterite, font=\scriptsize, anchor=south west]
    at (axis cs:2.1,7.4) {$x = 2$};
\end{axis}
\end{tikzpicture}
\end{center}

Les deux points marqués sont les extremums. On voit sur le dessin ce que le
tableau annonçait : la branche de gauche est entièrement sous l'asymptote
oblique, celle de droite entièrement au-dessus.
\end{exemple}

\begin{entrainement}[études de type 1]
\exo[\niveau{2}] $f(x) = \dfrac{x^{2}+1}{x}$ : mener les étapes 1 à 5 du plan
d'étude.

\exo[\niveau{2}] $f(x) = \dfrac{x^{2}-x-2}{x-1}$ : montrer que
$f(x) = x-\dfrac{2}{x-1}$, puis donner les asymptotes et la position de la
courbe par rapport à l'asymptote oblique.

\exo[\niveau{3}] $f(x) = \dfrac{2x^{2}-3x+1}{x-1}$ : simplifier $f(x)$ pour
$x \neq 1$ et en déduire la nature de la courbe. Que se passe-t-il en
$x = 1$ ?
\end{entrainement}

\begin{corrige}[études de type 1]
\textbf{1.} $\Df = \R^{*}$ ; $f(x) = x+\frac1x$ est impaire ;
$f' (x)= \frac{x^{2}-1}{x^{2}}$, positive à l'extérieur de
$\intervalle{-1}{1}$ ; maximum local $-2$ en $-1$, minimum local $2$ en
$1$ ; asymptotes $x = 0$ et $y = x$.

\textbf{2.} $x(x-1)-2 = x^{2}-x-2$ \checkmark. Asymptotes : $x = 1$ et
$y = x$. Le reste $-\frac{2}{x-1}$ est négatif pour $x>1$ : la courbe est
sous l'asymptote à droite de $1$, au-dessus à gauche.

\textbf{3.} $2x^{2}-3x+1 = (x-1)(2x-1)$, donc pour $x \neq 1$,
$f(x) = 2x-1$ : la courbe est une droite privée du point d'abscisse $1$. En
$x = 1$, la fonction n'est pas définie : le graphe présente un « trou » au
point $(1\,;\,1)$.
\end{corrige}

% ===========================================================================
\section{Deuxième type : \texorpdfstring{$f(x) = \dfrac{ax^{2}+bx+c}{dx^{2}+ex+l}$}{quotient de deux polynomes de degre 2}}

Ici le degré est le même en haut et en bas : la courbe admet une
\textbf{asymptote horizontale}, d'ordonnée le quotient des coefficients
dominants. Les valeurs interdites, elles, sont les racines du dénominateur —
il peut y en avoir deux, une, ou aucune.

\begin{exemple}[deux valeurs interdites]
Soit $f(x) = \dfrac{x^{2}+2}{x^{2}-1}$.

\emph{1.} $x^{2}-1 = 0$ pour $x = \pm1$, donc
$\Df = \R\setminus\{-1\,;\,1\}$.

\emph{2.} $\Df$ est symétrique et $f(-x) = f(x)$ : la fonction est
\textbf{paire}. Il suffira donc de l'étudier sur $\intervallefo{0}{+\infty}$
et de compléter le tracé par symétrie d'axe $(Oy)$ — une économie
considérable.

\emph{3.} $\limpinf f(x) = \limpinf \frac{x^{2}}{x^{2}} = 1$ : la droite
$y = 1$ est asymptote horizontale. En $1$, le numérateur tend vers $3>0$ et
le dénominateur vers $0$, négatif à gauche, positif à droite : $-\infty$ puis
$+\infty$. Les droites $x = -1$ et $x = 1$ sont asymptotes verticales.

\emph{4.}
\[
  f'(x) = \frac{2x\left(x^{2}-1\right)-\left(x^{2}+2\right)2x}
               {\left(x^{2}-1\right)^{2}}
        = \frac{-6x}{\left(x^{2}-1\right)^{2}} .
\]

\emph{5.} Le dénominateur est positif : $f'$ a le signe de $-6x$, donc $f$
croît sur les intervalles à gauche de $0$ et décroît à droite. En $0$,
$f(0) = -2$ : c'est un maximum local, sur l'intervalle
$\intervalleo{-1}{1}$.
\end{exemple}

\begin{avertissement}
Sur cet exemple, $f$ est croissante sur $\intervalleo{-\infty}{-1}$
\emph{et} sur $\intervalleo{-1}{0}$ — mais pas sur leur réunion : la
fonction saute de $+\infty$ à $-\infty$ en traversant $-1$. On ne réunit
jamais deux intervalles séparés par une valeur interdite.
\end{avertissement}

\begin{entrainement}[études de type 2]
\exo[\niveau{2}] $f(x) = \dfrac{2x^{2}}{x^{2}+1}$ : $\Df$, parité, limite en
$\pm\infty$, dérivée, variations.

\exo[\niveau{2}] $f(x) = \dfrac{x^{2}-4}{x^{2}-9}$ : $\Df$, asymptotes,
intersections avec les axes.

\exo[\niveau{3}] $f(x) = \dfrac{3x^{2}+1}{x^{2}-4}$ : mener l'étude
complète, tracé compris.
\end{entrainement}

\begin{corrige}[études de type 2]
\textbf{1.} $\Df = \R$ (le dénominateur ne s'annule jamais) ; fonction
paire ; $\limpinf f = 2$, asymptote $y = 2$ ;
$f'(x) = \frac{4x}{\left(x^{2}+1\right)^{2}}$, du signe de $x$ : $f$ décroît
sur $\R^{-}$, croît sur $\R^{+}$, minimum $f(0) = 0$.

\textbf{2.} $\Df = \R\setminus\{-3\,;\,3\}$ ; asymptotes $x = -3$, $x = 3$ et
$y = 1$ ; la courbe coupe $(Ox)$ en $-2$ et $2$, et $(Oy)$ en
$\left(0\,;\,\frac49\right)$.

\textbf{3.} $\Df = \R\setminus\{-2\,;\,2\}$ ; paire ; asymptote horizontale
$y = 3$ et asymptotes verticales $x = \pm2$ ;
$f'(x) = \frac{-26x}{\left(x^{2}-4\right)^{2}}$, du signe de $-x$ ; maximum
local $f(0) = -\frac14$ sur $\intervalleo{-2}{2}$.
\end{corrige}

% ===========================================================================
\section{Position d'une courbe par rapport à une droite}

\begin{methode}{comparer une courbe et une droite}
Pour situer $\Cf$ par rapport à la droite $(D) : y = ax+b$ :
\begin{enumerate}[label=\textbf{\arabic*.}, leftmargin=*, itemsep=0.15em]
  \item calculer la différence $d(x) = f(x)-(ax+b)$ ;
  \item réduire au même dénominateur et factoriser ;
  \item étudier le \textbf{signe} de $d(x)$ dans un tableau ;
  \item conclure : $d(x) > 0$ signifie « la courbe est au-dessus »,
        $d(x) < 0$ « en dessous », $d(x) = 0$ donne les points
        d'intersection.
\end{enumerate}
\end{methode}

\begin{exemple}[une courbe qui traverse sa droite]
Soit $f(x) = \dfrac{x^{3}}{x^{2}+1}$ et $(D) : y = x$. Alors
\[
  f(x)-x = \frac{x^{3}-x\left(x^{2}+1\right)}{x^{2}+1} = \frac{-x}{x^{2}+1},
\]
qui est du signe de $-x$. La courbe est donc au-dessus de $(D)$ pour
$x<0$, en dessous pour $x>0$, et les deux se coupent en l'origine : une
courbe traverse parfaitement son asymptote, et cela n'a rien d'anormal.
\end{exemple}

\begin{entrainement}[positions relatives]
\exo[\niveau{1}] $f(x) = x+\dfrac{3}{x}$ et $(D) : y = x$ : étudier la
position.

\exo[\niveau{2}] $f(x) = \dfrac{2x^{2}-1}{x}$ et $(D) : y = 2x$ : étudier la
position et donner les points d'intersection s'il en existe.

\exo[\niveau{3}] $f(x) = \dfrac{x^{2}+x+1}{x+1}$ et $(D) : y = x$ : montrer
que la courbe est toujours au-dessus de $(D)$ sur
$\intervalleo{-1}{+\infty}$.
\end{entrainement}

\begin{corrige}[positions relatives]
\textbf{1.} $f(x)-x = \frac3x$, du signe de $x$ : au-dessus pour $x>0$, en
dessous pour $x<0$ ; aucun point d'intersection.
\textbf{2.} $f(x)-2x = -\frac1x$, du signe de $-x$ : au-dessus pour $x<0$,
en dessous pour $x>0$ ; aucun point d'intersection, la différence ne
s'annulant jamais.
\textbf{3.} $f(x)-x = \frac{x^{2}+x+1-x(x+1)}{x+1} = \frac{1}{x+1}$, qui est
strictement positif dès que $x>-1$.
\end{corrige}

\begin{qcm}
\exo Si $\limpinf f(x) = 2$ et $\limminf f(x) = 2$, la courbe admet :
\textbf{a)} deux asymptotes horizontales distinctes \quad
\textbf{b)} une seule asymptote horizontale, $y = 2$ \quad
\textbf{c)} une asymptote oblique \quad \textbf{d)} aucune asymptote
\reponseqcm{b}

\exo Une fonction paire dont on connaît l'étude sur $\intervallefo{0}{+\infty}$ :
\textbf{a)} doit être réétudiée sur $\R^{-}$ \quad
\textbf{b)} se complète par symétrie d'axe $(Oy)$ \quad
\textbf{c)} se complète par symétrie de centre $O$ \quad
\textbf{d)} n'est pas définie sur $\R^{-}$ \reponseqcm{b}

\exo La courbe de $f$ coupe l'axe des abscisses aux points d'abscisses :
\textbf{a)} les valeurs interdites \quad \textbf{b)} les solutions de
$f(x) = 0$ \quad \textbf{c)} les solutions de $f'(x) = 0$ \quad
\textbf{d)} $f(0)$ \reponseqcm{b}

\exo Pour $f(x) = \dfrac{ax^{2}+bx+c}{dx+e}$ avec $a \neq 0$ et $d \neq 0$,
la courbe admet en $+\infty$ :
\textbf{a)} une asymptote horizontale \quad \textbf{b)} une asymptote
oblique \quad \textbf{c)} une asymptote verticale \quad \textbf{d)} aucune
asymptote \reponseqcm{b}
\end{qcm}

% ===========================================================================
\begin{annales}
Les études complètes des sujets malgaches reposent sur le logarithme ou
l'exponentielle et figurent donc aux chapitres 7 et 8. Les énoncés ci-dessous
sont des études réelles de fonctions rationnelles, du même fonds vérifié.

\exoann{Baccalauréat blanc, Lycée Philibert Tsiranana, session 2022-2023 —
problème}
Soit $f(x) = \dfrac{x^{2}-3x+4}{x-2}$, de courbe $\Cf$ dans un repère
orthonormé d'unité $1$ cm.
\begin{enumerate}[label=\textbf{\alph*)}, leftmargin=*, itemsep=0.15em]
  \item Déterminer $\Df$ et les limites de $f$ aux bornes de $\Df$.
  \item Vérifier que $f(x) = x-1+\dfrac{2}{x-2}$ et en déduire l'asymptote
        oblique $(D)$.
  \item Étudier la position de $\Cf$ par rapport à $(D)$.
  \item Montrer que $f'(x) = \dfrac{x^{2}-4x+2}{(x-2)^{2}}$, puis dresser le
        tableau de variation.
  \item Tracer $(D)$, l'asymptote verticale et $\Cf$.
\end{enumerate}

\exoann{Baccalauréat probatoire, Cameroun, séries D et TI, session 2023 —
exercice 4}
Soit $f(x) = x+\dfrac{1}{x+1}$.
\begin{enumerate}[label=\textbf{\alph*)}, leftmargin=*, itemsep=0.15em]
  \item Déterminer $\Df$ et les limites aux bornes.
  \item Montrer que la droite $y = x$ et la droite $x = -1$ sont asymptotes
        à la courbe.
  \item Montrer que $f'(x) = \dfrac{x(x+2)}{(x+1)^{2}}$ et dresser le
        tableau de variation.
  \item Étudier la position de la courbe par rapport à la droite $y = x$.
  \item Tracer la courbe et ses asymptotes.
\end{enumerate}

\exoann{Baccalauréat probatoire, Cameroun, série D, session 2022 —
exercice 4}
Soit $f(x) = \dfrac{3x}{3+4x}$. Déterminer $\Df$, les limites aux bornes, la
dérivée, le tableau de variation, puis tracer la courbe avec ses deux
asymptotes.
\end{annales}

\begin{corrige}[annales]
\textbf{1.} C'est exactement l'étude menée en exemple dans ce chapitre :
$\Df = \R\setminus\{2\}$ ; limites $\pm\infty$ en $\pm\infty$, $-\infty$ en
$2^{-}$ et $+\infty$ en $2^{+}$ ; asymptote oblique $y = x-1$, courbe
au-dessus pour $x>2$ et en dessous pour $x<2$ ; extremums
$f(2-\sqrt2) = 1-2\sqrt2 \approx -1{,}83$ (maximum local) et
$f(2+\sqrt2) = 1+2\sqrt2 \approx 4{,}83$ (minimum local).

\textbf{2.} $\Df = \R\setminus\{-1\}$ ; $f(x)-x = \frac{1}{x+1} \to 0$, donc
$y = x$ est asymptote oblique, et la courbe est au-dessus de cette droite
pour $x>-1$, en dessous pour $x<-1$ ; $x = -1$ est asymptote verticale
($-\infty$ à gauche, $+\infty$ à droite) ; variations : croissante jusqu'en
$-2$ (maximum $-3$), décroissante jusqu'en $-1$, décroissante de $-1$ à $0$
(minimum $1$), puis croissante.

\textbf{3.} $\Df = \R\setminus\left\{-\frac34\right\}$ ; asymptotes
$y = \frac34$ et $x = -\frac34$ ; $f'(x) = \frac{9}{(3+4x)^{2}} > 0$ : la
fonction est strictement croissante sur chacun des deux intervalles. La
courbe coupe les deux axes à l'origine.
\end{corrige}

% ===========================================================================
\begin{jecherche}[reconstituer une fonction à partir de sa courbe]
La courbe ci-dessous est celle d'une fonction de la forme
$f(x) = \dfrac{ax+b}{x+c}$.

\begin{center}
\begin{tikzpicture}
\begin{axis}[
  width = 9.6cm, height = 5.6cm,
  axis lines = middle, xlabel = {$x$}, ylabel = {$y$},
  xmin = -5.4, xmax = 5.4, ymin = -4.4, ymax = 6.4,
  xtick = {-5,-4,-3,-2,-1,1,2,3,4,5}, ytick = {-4,-2,2,4,6},
  axis line style = {bmtGris},
  tick label style = {font=\scriptsize, color=bmtGris},
  grid = both, grid style = {bmtGrisClair, line width=0.3pt},
  clip = true, restrict y to domain = -10:12,
]
  \addplot[bmtIndigo, line width=1.4pt, domain=-5.3:0.88, samples=200]
    {(2*x+4)/(x-1)};
  \addplot[bmtIndigo, line width=1.4pt, domain=1.14:5.3, samples=200]
    {(2*x+4)/(x-1)};
  \addplot[bmtLaterite, line width=0.9pt, dash pattern=on 3pt off 2pt,
           domain=-5.3:5.3] {2};
  \draw[bmtLaterite, line width=0.9pt, dash pattern=on 3pt off 2pt]
    (axis cs:1,-4.3) -- (axis cs:1,6.3);
\end{axis}
\end{tikzpicture}
\end{center}

\begin{enumerate}[label=\textbf{\arabic*.}, leftmargin=*, itemsep=0.3em]
  \item Lire l'asymptote verticale et en déduire $c$.
  \item Lire l'asymptote horizontale et en déduire $a$.
  \item La courbe coupe l'axe des abscisses en $-2$. En déduire $b$.
  \item Vérifier votre réponse en calculant $f(0)$ et en le comparant au
        graphique.
  \item Calculer $f'(x)$ et vérifier que le signe obtenu correspond bien au
        sens de variation lu sur la figure.
\end{enumerate}
\end{jecherche}

\begin{corrige}[reconstituer une fonction à partir de sa courbe]
\textbf{1.} L'asymptote verticale est $x = 1$, donc $x+c$ s'annule en $1$ :
$c = -1$.

\textbf{2.} L'asymptote horizontale est $y = 2$, et
$\limpinf \frac{ax+b}{x-1} = a$ : donc $a = 2$.

\textbf{3.} $f(-2) = 0$ impose $2\times(-2)+b = 0$, soit $b = 4$. Ainsi
$f(x) = \frac{2x+4}{x-1}$.

\textbf{4.} $f(0) = \frac{4}{-1} = -4$, ce que confirme le graphique.

\textbf{5.} $f'(x) = \frac{2(x-1)-(2x+4)}{(x-1)^{2}} = \frac{-6}{(x-1)^{2}}
< 0$ : la fonction décroît sur chacun des deux intervalles, ce qui est bien
ce que montre la figure.
\end{corrige}

% ===========================================================================
\begin{vraifaux}
\exo Une fonction rationnelle dont le numérateur et le dénominateur ont le
même degré admet une asymptote horizontale.

\exo Si une courbe admet une asymptote oblique, elle ne peut pas avoir
d'asymptote verticale.

\exo Une fonction paire ne peut pas avoir d'asymptote oblique.

\exo Les points d'intersection de $\Cf$ avec l'axe des abscisses sont donnés
par $f'(x) = 0$.

\exo Si $f(x)-(x+1)$ garde un signe constant, la courbe ne coupe jamais la
droite $y = x+1$.
\end{vraifaux}

\begin{corrige}[vrai ou faux]
\textbf{1.} Vrai — d'ordonnée le quotient des coefficients dominants.
\textbf{2.} Faux — l'exemple du chapitre en a les deux.
\textbf{3.} Faux, mais presque : si elle en a une en $+\infty$, la symétrie
en donne une autre en $-\infty$, de coefficient directeur opposé. Une même
droite ne peut pas être asymptote des deux côtés.
\textbf{4.} Faux — c'est $f(x) = 0$ ; $f'(x) = 0$ donne les tangentes
horizontales.
\textbf{5.} Vrai — couper la droite exigerait que la différence s'annule.
\end{corrige}

% ===========================================================================
\begin{jemeteste}
Trente minutes.

Soit $f(x) = \dfrac{x^{2}-4}{x-1}$.

\begin{enumerate}[label=\textbf{\arabic*.}, leftmargin=*, itemsep=0.3em]
  \item Déterminer $\Df$.
  \item Vérifier que $f(x) = x+1-\dfrac{3}{x-1}$.
  \item En déduire les trois asymptotes ou limites remarquables.
  \item Étudier la position de $\Cf$ par rapport à la droite $y = x+1$.
  \item Calculer $f'(x)$ et donner son signe.
  \item Dresser le tableau de variation.
\end{enumerate}
\end{jemeteste}

\begin{corrige}[je me teste]
\textbf{1.} $\Df = \R\setminus\{1\}$.
\textbf{2.} $(x+1)(x-1)-3 = x^{2}-1-3 = x^{2}-4$. \checkmark
\textbf{3.} $y = x+1$ est asymptote oblique en $\pm\infty$ ; $x = 1$ est
asymptote verticale, avec $+\infty$ en $1^{-}$ et $-\infty$ en $1^{+}$ (le
signe vient du $-3$ au numérateur du reste).
\textbf{4.} $f(x)-(x+1) = -\frac{3}{x-1}$ : positif pour $x<1$ (courbe
au-dessus), négatif pour $x>1$ (courbe en dessous).
\textbf{5.} $f'(x) = 1+\frac{3}{(x-1)^{2}} > 0$ : toujours strictement
positive.
\textbf{6.} $f$ est strictement croissante sur $\intervalleo{-\infty}{1}$
(de $-\infty$ à $+\infty$) et sur $\intervalleo{1}{+\infty}$ (de $-\infty$
à $+\infty$).
\end{corrige}

% ===========================================================================
\begin{commeaubac}
\bmtsource{Exercice original au format de l'épreuve — série A \textbullet\ 10 points, format du problème}
Soit $f$ la fonction définie sur $\R\setminus\{-1\}$ par
\[
  f(x) = \frac{x^{2}+2x+2}{x+1},
\]
de courbe $\Cf$ dans un repère orthonormé d'unité $1$ cm.

\begin{enumerate}[label=\textbf{\arabic*.}, leftmargin=*, itemsep=0.25em]
  \item Justifier que $\Df = \R\setminus\{-1\}$. \bareme{0,5}
  \item Vérifier que, pour tout $x \neq -1$,
        $f(x) = x+1+\dfrac{1}{x+1}$. \bareme{1}
  \item En déduire :
    \begin{enumerate}[label=\textbf{\alph*)}, leftmargin=1.4em, itemsep=0.15em]
      \item les limites de $f$ en $-\infty$ et en $+\infty$ ;
            \bareme{1}
      \item les limites de $f$ en $-1^{-}$ et en $-1^{+}$, et l'asymptote
            verticale qui en résulte. \bareme{1}
    \end{enumerate}
  \item Montrer que la droite $(D) : y = x+1$ est asymptote oblique à $\Cf$
        en $+\infty$ et en $-\infty$, puis étudier la position de $\Cf$ par
        rapport à $(D)$. \bareme{1,5}
  \item Montrer que $f'(x) = \dfrac{x(x+2)}{(x+1)^{2}}$, étudier son signe et
        dresser le tableau de variation de $f$. \bareme{2,5}
  \item Écrire une équation de la tangente $(T)$ à $\Cf$ au point d'abscisse
        $0$. \bareme{1}
  \item Tracer $(D)$, l'asymptote verticale, $(T)$ et $\Cf$. \bareme{1,5}
\end{enumerate}
\end{commeaubac}

\begin{corrige}[comme au baccalauréat]
\textbf{1.} Le dénominateur $x+1$ s'annule pour $x = -1$ seulement.

\textbf{2.} $(x+1)(x+1)+1 = x^{2}+2x+1+1 = x^{2}+2x+2$. \checkmark

\textbf{3. a)} $\frac{1}{x+1} \to 0$, donc $f(x) \sim x+1$ :
$f \to -\infty$ en $-\infty$ et $+\infty$ en $+\infty$.
\textbf{b)} En $-1$, le terme $x+1$ tend vers $0$ et $\frac{1}{x+1}$ vers
$-\infty$ à gauche, $+\infty$ à droite. La droite $x = -1$ est asymptote
verticale.

\textbf{4.} $f(x)-(x+1) = \frac{1}{x+1} \to 0$ : $(D)$ est asymptote
oblique. Ce reste a le signe de $x+1$ : la courbe est au-dessus de $(D)$
pour $x>-1$, en dessous pour $x<-1$.

\textbf{5.} $f'(x) = 1-\frac{1}{(x+1)^{2}} = \frac{(x+1)^{2}-1}{(x+1)^{2}}
= \frac{x^{2}+2x}{(x+1)^{2}} = \frac{x(x+2)}{(x+1)^{2}}$. Le dénominateur
est positif ; $x(x+2)$ est positif à l'extérieur de $\intervalle{-2}{0}$.
Donc $f$ croît sur $\intervalleof{-\infty}{-2}$, décroît sur
$\intervallefo{-2}{-1}$ et sur $\intervalleof{-1}{0}$, croît sur
$\intervallefo{0}{+\infty}$. Maximum local $f(-2) = -2$, minimum local
$f(0) = 2$.

\textbf{6.} $f(0) = 2$ et $f'(0) = 0$ : la tangente est horizontale,
$(T) : y = 2$.

\textbf{7.} On trace d'abord les deux asymptotes en pointillé, puis la
tangente horizontale $y = 2$, puis la courbe : une branche en haut à droite
au-dessus de $(D)$, une branche en bas à gauche en dessous.
\end{corrige}

\begin{notepourlaclasse}
Ce chapitre ne s'enseigne pas, il se \emph{fait faire}. Une seule étude
menée par le professeur au tableau, du début au tracé, puis trois études
conduites par les élèves eux-mêmes, en autonomie, avec la liste des sept
étapes affichée au mur : c'est le rythme qui fonctionne.

Exigez le tracé sur papier millimétré et une unité respectée. Les sujets
précisent toujours l'unité graphique — $1$ cm, $2$ cm — et le tracé compte
pour deux points, autant qu'un calcul de dérivée.

Les élèves confondent régulièrement « position par rapport à l'asymptote »
et « signe de la fonction ». Faites-leur écrire la différence $f(x)-(ax+b)$
en toutes lettres avant tout calcul : le geste suffit à lever la confusion.
\end{notepourlaclasse}
